% Options for packages loaded elsewhere
% Options for packages loaded elsewhere
\PassOptionsToPackage{unicode}{hyperref}
\PassOptionsToPackage{hyphens}{url}
\PassOptionsToPackage{dvipsnames,svgnames,x11names}{xcolor}
%
\documentclass[
  ngerman,
  letterpaper,
  DIV=11]{scrreprt}
\usepackage{xcolor}
\usepackage{amsmath,amssymb}
\setcounter{secnumdepth}{-\maxdimen} % remove section numbering
\usepackage{iftex}
\ifPDFTeX
  \usepackage[T1]{fontenc}
  \usepackage[utf8]{inputenc}
  \usepackage{textcomp} % provide euro and other symbols
\else % if luatex or xetex
  \usepackage{unicode-math} % this also loads fontspec
  \defaultfontfeatures{Scale=MatchLowercase}
  \defaultfontfeatures[\rmfamily]{Ligatures=TeX,Scale=1}
\fi
\usepackage{lmodern}
\ifPDFTeX\else
  % xetex/luatex font selection
\fi
% Use upquote if available, for straight quotes in verbatim environments
\IfFileExists{upquote.sty}{\usepackage{upquote}}{}
\IfFileExists{microtype.sty}{% use microtype if available
  \usepackage[]{microtype}
  \UseMicrotypeSet[protrusion]{basicmath} % disable protrusion for tt fonts
}{}
\makeatletter
\@ifundefined{KOMAClassName}{% if non-KOMA class
  \IfFileExists{parskip.sty}{%
    \usepackage{parskip}
  }{% else
    \setlength{\parindent}{0pt}
    \setlength{\parskip}{6pt plus 2pt minus 1pt}}
}{% if KOMA class
  \KOMAoptions{parskip=half}}
\makeatother
% Make \paragraph and \subparagraph free-standing
\makeatletter
\ifx\paragraph\undefined\else
  \let\oldparagraph\paragraph
  \renewcommand{\paragraph}{
    \@ifstar
      \xxxParagraphStar
      \xxxParagraphNoStar
  }
  \newcommand{\xxxParagraphStar}[1]{\oldparagraph*{#1}\mbox{}}
  \newcommand{\xxxParagraphNoStar}[1]{\oldparagraph{#1}\mbox{}}
\fi
\ifx\subparagraph\undefined\else
  \let\oldsubparagraph\subparagraph
  \renewcommand{\subparagraph}{
    \@ifstar
      \xxxSubParagraphStar
      \xxxSubParagraphNoStar
  }
  \newcommand{\xxxSubParagraphStar}[1]{\oldsubparagraph*{#1}\mbox{}}
  \newcommand{\xxxSubParagraphNoStar}[1]{\oldsubparagraph{#1}\mbox{}}
\fi
\makeatother

\usepackage{color}
\usepackage{fancyvrb}
\newcommand{\VerbBar}{|}
\newcommand{\VERB}{\Verb[commandchars=\\\{\}]}
\DefineVerbatimEnvironment{Highlighting}{Verbatim}{commandchars=\\\{\}}
% Add ',fontsize=\small' for more characters per line
\usepackage{framed}
\definecolor{shadecolor}{RGB}{241,243,245}
\newenvironment{Shaded}{\begin{snugshade}}{\end{snugshade}}
\newcommand{\AlertTok}[1]{\textcolor[rgb]{0.68,0.00,0.00}{#1}}
\newcommand{\AnnotationTok}[1]{\textcolor[rgb]{0.37,0.37,0.37}{#1}}
\newcommand{\AttributeTok}[1]{\textcolor[rgb]{0.40,0.45,0.13}{#1}}
\newcommand{\BaseNTok}[1]{\textcolor[rgb]{0.68,0.00,0.00}{#1}}
\newcommand{\BuiltInTok}[1]{\textcolor[rgb]{0.00,0.23,0.31}{#1}}
\newcommand{\CharTok}[1]{\textcolor[rgb]{0.13,0.47,0.30}{#1}}
\newcommand{\CommentTok}[1]{\textcolor[rgb]{0.37,0.37,0.37}{#1}}
\newcommand{\CommentVarTok}[1]{\textcolor[rgb]{0.37,0.37,0.37}{\textit{#1}}}
\newcommand{\ConstantTok}[1]{\textcolor[rgb]{0.56,0.35,0.01}{#1}}
\newcommand{\ControlFlowTok}[1]{\textcolor[rgb]{0.00,0.23,0.31}{\textbf{#1}}}
\newcommand{\DataTypeTok}[1]{\textcolor[rgb]{0.68,0.00,0.00}{#1}}
\newcommand{\DecValTok}[1]{\textcolor[rgb]{0.68,0.00,0.00}{#1}}
\newcommand{\DocumentationTok}[1]{\textcolor[rgb]{0.37,0.37,0.37}{\textit{#1}}}
\newcommand{\ErrorTok}[1]{\textcolor[rgb]{0.68,0.00,0.00}{#1}}
\newcommand{\ExtensionTok}[1]{\textcolor[rgb]{0.00,0.23,0.31}{#1}}
\newcommand{\FloatTok}[1]{\textcolor[rgb]{0.68,0.00,0.00}{#1}}
\newcommand{\FunctionTok}[1]{\textcolor[rgb]{0.28,0.35,0.67}{#1}}
\newcommand{\ImportTok}[1]{\textcolor[rgb]{0.00,0.46,0.62}{#1}}
\newcommand{\InformationTok}[1]{\textcolor[rgb]{0.37,0.37,0.37}{#1}}
\newcommand{\KeywordTok}[1]{\textcolor[rgb]{0.00,0.23,0.31}{\textbf{#1}}}
\newcommand{\NormalTok}[1]{\textcolor[rgb]{0.00,0.23,0.31}{#1}}
\newcommand{\OperatorTok}[1]{\textcolor[rgb]{0.37,0.37,0.37}{#1}}
\newcommand{\OtherTok}[1]{\textcolor[rgb]{0.00,0.23,0.31}{#1}}
\newcommand{\PreprocessorTok}[1]{\textcolor[rgb]{0.68,0.00,0.00}{#1}}
\newcommand{\RegionMarkerTok}[1]{\textcolor[rgb]{0.00,0.23,0.31}{#1}}
\newcommand{\SpecialCharTok}[1]{\textcolor[rgb]{0.37,0.37,0.37}{#1}}
\newcommand{\SpecialStringTok}[1]{\textcolor[rgb]{0.13,0.47,0.30}{#1}}
\newcommand{\StringTok}[1]{\textcolor[rgb]{0.13,0.47,0.30}{#1}}
\newcommand{\VariableTok}[1]{\textcolor[rgb]{0.07,0.07,0.07}{#1}}
\newcommand{\VerbatimStringTok}[1]{\textcolor[rgb]{0.13,0.47,0.30}{#1}}
\newcommand{\WarningTok}[1]{\textcolor[rgb]{0.37,0.37,0.37}{\textit{#1}}}

\usepackage{longtable,booktabs,array}
\usepackage{calc} % for calculating minipage widths
% Correct order of tables after \paragraph or \subparagraph
\usepackage{etoolbox}
\makeatletter
\patchcmd\longtable{\par}{\if@noskipsec\mbox{}\fi\par}{}{}
\makeatother
% Allow footnotes in longtable head/foot
\IfFileExists{footnotehyper.sty}{\usepackage{footnotehyper}}{\usepackage{footnote}}
\makesavenoteenv{longtable}
\usepackage{graphicx}
\makeatletter
\newsavebox\pandoc@box
\newcommand*\pandocbounded[1]{% scales image to fit in text height/width
  \sbox\pandoc@box{#1}%
  \Gscale@div\@tempa{\textheight}{\dimexpr\ht\pandoc@box+\dp\pandoc@box\relax}%
  \Gscale@div\@tempb{\linewidth}{\wd\pandoc@box}%
  \ifdim\@tempb\p@<\@tempa\p@\let\@tempa\@tempb\fi% select the smaller of both
  \ifdim\@tempa\p@<\p@\scalebox{\@tempa}{\usebox\pandoc@box}%
  \else\usebox{\pandoc@box}%
  \fi%
}
% Set default figure placement to htbp
\def\fps@figure{htbp}
\makeatother



\ifLuaTeX
\usepackage[bidi=basic]{babel}
\else
\usepackage[bidi=default]{babel}
\fi
% get rid of language-specific shorthands (see #6817):
\let\LanguageShortHands\languageshorthands
\def\languageshorthands#1{}
\ifLuaTeX
  \usepackage[german]{selnolig} % disable illegal ligatures
\fi


\setlength{\emergencystretch}{3em} % prevent overfull lines

\providecommand{\tightlist}{%
  \setlength{\itemsep}{0pt}\setlength{\parskip}{0pt}}



 


\usepackage[top=2cm, bottom=2.5cm]{geometry}
\usepackage{enumitem}
\usepackage{changepage}
\setlist[description]{style=nextline, leftmargin=2.5em}
\newcommand{\cmd}[1]{\texttt{#1}\xspace}
\KOMAoption{captions}{tableheading}
\makeatletter
\@ifpackageloaded{caption}{}{\usepackage{caption}}
\AtBeginDocument{%
\ifdefined\contentsname
  \renewcommand*\contentsname{Inhaltsverzeichnis}
\else
  \newcommand\contentsname{Inhaltsverzeichnis}
\fi
\ifdefined\listfigurename
  \renewcommand*\listfigurename{Abbildungsverzeichnis}
\else
  \newcommand\listfigurename{Abbildungsverzeichnis}
\fi
\ifdefined\listtablename
  \renewcommand*\listtablename{Tabellenverzeichnis}
\else
  \newcommand\listtablename{Tabellenverzeichnis}
\fi
\ifdefined\figurename
  \renewcommand*\figurename{Abbildung}
\else
  \newcommand\figurename{Abbildung}
\fi
\ifdefined\tablename
  \renewcommand*\tablename{Tabelle}
\else
  \newcommand\tablename{Tabelle}
\fi
}
\@ifpackageloaded{float}{}{\usepackage{float}}
\floatstyle{ruled}
\@ifundefined{c@chapter}{\newfloat{codelisting}{h}{lop}}{\newfloat{codelisting}{h}{lop}[chapter]}
\floatname{codelisting}{Listing}
\newcommand*\listoflistings{\listof{codelisting}{Listingverzeichnis}}
\makeatother
\makeatletter
\makeatother
\makeatletter
\@ifpackageloaded{caption}{}{\usepackage{caption}}
\@ifpackageloaded{subcaption}{}{\usepackage{subcaption}}
\makeatother
\usepackage{bookmark}
\IfFileExists{xurl.sty}{\usepackage{xurl}}{} % add URL line breaks if available
\urlstyle{same}
\hypersetup{
  pdftitle={Git-HTTP kennwortlos},
  pdfauthor={Wolfgang Höfer},
  pdflang={de},
  colorlinks=true,
  linkcolor={blue},
  filecolor={Maroon},
  citecolor={Blue},
  urlcolor={Blue},
  pdfcreator={LaTeX via pandoc}}


\title{Git-HTTP kennwortlos}
\author{Wolfgang Höfer}
\date{}
\begin{document}
\maketitle


\chapter[Anleitung]{\texorpdfstring{Anleitung\footnote{Dieses Werk steht
  unter der \textbf{Creative Commons Attribution-ShareAlike 4.0
  International License (CC BY-SA 4.0)}.}}{Anleitung}}\label{anleitungcc}

Eine unsichere, aber schnelle Variante ist ein Git-Server, der http ohne
Kennwort unterstützt.

Nachfolgend die relevanten Dateien:

\begin{itemize}
\tightlist
\item
  docker-compose.yml
\item
  nginx.conf
\item
  Dockerfile
\item
  startup.sh
\end{itemize}

Der Start und Einsatz des Containers wird danach beschrieben

\textbf{docker-compose.yml}

\begin{Shaded}
\begin{Highlighting}[]
\ExtensionTok{version:} \StringTok{\textquotesingle{}3.8\textquotesingle{}}

\ExtensionTok{services:}
  \ExtensionTok{git{-}nginx:}
    \ExtensionTok{build:}
      \ExtensionTok{context:}\NormalTok{ .}
      \ExtensionTok{dockerfile:}\NormalTok{ Dockerfile}
    \ExtensionTok{restart:}\NormalTok{ always}
    \ExtensionTok{container\_name:}\NormalTok{ git{-}nginx{-}container}
    \ExtensionTok{ports:}
      \ExtensionTok{{-}} \StringTok{"8080:80"}     \CommentTok{\# HTTP Port}
    \ExtensionTok{volumes:}
      \ExtensionTok{{-}}\NormalTok{ ./repo:/git}
      \ExtensionTok{{-}}\NormalTok{ ./repositories:/var/www/git/repositories }\CommentTok{\# nicht klar wofür}
      \ExtensionTok{{-}}\NormalTok{ ./nginx.conf:/etc/nginx/nginx.conf  }
\end{Highlighting}
\end{Shaded}

\textbf{nginx.conf}

\begin{Shaded}
\begin{Highlighting}[]
\ExtensionTok{worker\_processes}\NormalTok{  1}\KeywordTok{;}

\ExtensionTok{error\_log}\NormalTok{  /var/log/nginx/error.log}\KeywordTok{;}
\ExtensionTok{pid}\NormalTok{ /run/nginx.pid}\KeywordTok{;}
\ExtensionTok{user}\NormalTok{ root}\KeywordTok{;}

\ExtensionTok{events}\NormalTok{ \{}
    \ExtensionTok{worker\_connections}\NormalTok{  1024}\KeywordTok{;}
\ErrorTok{\}}

\ExtensionTok{http}\NormalTok{ \{}
    \ExtensionTok{server}\NormalTok{ \{}
        \ExtensionTok{listen}  \PreprocessorTok{*}\NormalTok{:80}\KeywordTok{;}

        \ExtensionTok{root}\NormalTok{ /www/empty/}\KeywordTok{;}
        \ExtensionTok{index}\NormalTok{ index.html}\KeywordTok{;}

        \ExtensionTok{server\_name} \VariableTok{$hostname}\KeywordTok{;}
        \ExtensionTok{access\_log}\NormalTok{ /var/log/nginx/access.log}\KeywordTok{;}

        \CommentTok{\#error\_page 404 /404.html;}

        \CommentTok{\#auth\_basic            "Restricted";}
        \CommentTok{\#auth\_basic\_user\_file  /www/htpasswd;}

        \ExtensionTok{location}\NormalTok{ \textasciitilde{} /git}\ErrorTok{(}\ExtensionTok{/.*}\KeywordTok{)} \KeywordTok{\{}
            \CommentTok{\# Set chunks to unlimited, as the bodies can be huge}
            \ExtensionTok{client\_max\_body\_size}\NormalTok{            0}\KeywordTok{;}

            \ExtensionTok{fastcgi\_param}\NormalTok{ SCRIPT\_FILENAME /usr/libexec/git{-}core/git{-}http{-}backend}\KeywordTok{;}
            \ExtensionTok{include}\NormalTok{ fastcgi\_params}\KeywordTok{;}
            \ExtensionTok{fastcgi\_param}\NormalTok{ GIT\_HTTP\_EXPORT\_ALL }\StringTok{""}\KeywordTok{;}
            \ExtensionTok{fastcgi\_param}\NormalTok{ GIT\_PROJECT\_ROOT /git}\KeywordTok{;}
            \ExtensionTok{fastcgi\_param}\NormalTok{ PATH\_INFO }\VariableTok{$1}\KeywordTok{;}

            \ExtensionTok{fastcgi\_param}\NormalTok{ REQUEST\_METHOD }\VariableTok{$request\_method}\KeywordTok{;}
            \ExtensionTok{fastcgi\_param}\NormalTok{ QUERY\_STRING }\VariableTok{$query\_string}\KeywordTok{;}

            \CommentTok{\# Forward REMOTE\_USER as we want to know when we are authenticated}
            \ExtensionTok{fastcgi\_param}\NormalTok{   REMOTE\_USER     }\VariableTok{$remote\_user}\KeywordTok{;}
            \ExtensionTok{fastcgi\_pass}\NormalTok{    unix:/run/fcgi.sock}\KeywordTok{;}
        \KeywordTok{\}}
    \ErrorTok{\}}
\end{Highlighting}
\end{Shaded}

\textbf{Dockerfile}

\begin{Shaded}
\begin{Highlighting}[]
\ExtensionTok{FROM}\NormalTok{ alpine:latest}

\ExtensionTok{MAINTAINER}\NormalTok{ Anthony Hogg anthony@hogg.fr}

\CommentTok{\# The container listens on port 80, map as needed}
\ExtensionTok{EXPOSE}\NormalTok{ 80}

\CommentTok{\# Der Ort im Image, wo die Repos liegen werden}
\CommentTok{\# Da wir ein benanntes Volume wollen, überschreiben wir }
\CommentTok{\# das im compose{-}File!}
\CommentTok{\# VOLUME ["/git"]}

\CommentTok{\# We need the following:}
\CommentTok{\# {-} git{-}daemon, because that gets us the git{-}http{-}backend CGI script}
\CommentTok{\# {-} fcgiwrap, because that is how nginx does CGI}
\CommentTok{\# {-} spawn{-}fcgi, to launch fcgiwrap and to create the unix socket}
\CommentTok{\# {-} nginx, because it is our frontend}
\ExtensionTok{RUN}\NormalTok{ apk add }\AttributeTok{{-}{-}update}\NormalTok{ nginx }\KeywordTok{\&\&} \DataTypeTok{\textbackslash{}}
    \ExtensionTok{apk}\NormalTok{ add }\AttributeTok{{-}{-}update}\NormalTok{ git{-}daemon }\KeywordTok{\&\&} \DataTypeTok{\textbackslash{}}
    \ExtensionTok{apk}\NormalTok{ add }\AttributeTok{{-}{-}update}\NormalTok{ fcgiwrap }\KeywordTok{\&\&} \DataTypeTok{\textbackslash{}}
    \ExtensionTok{apk}\NormalTok{ add }\AttributeTok{{-}{-}update}\NormalTok{ spawn{-}fcgi }\KeywordTok{\&\&} \DataTypeTok{\textbackslash{}}
    \FunctionTok{rm} \AttributeTok{{-}rf}\NormalTok{ /var/cache/apk/}\PreprocessorTok{*}

\ExtensionTok{COPY}\NormalTok{ nginx.conf /etc/nginx/nginx.conf}

\CommentTok{\# launch fcgiwrap via spawn{-}fcgi; launch nginx in the foreground}
\CommentTok{\# so the container doesn\textquotesingle{}t die on us; supposedly we should be}
\CommentTok{\# using supervisord or something like that instead, but this}
\CommentTok{\# will do}
\ExtensionTok{CMD}\NormalTok{ spawn{-}fcgi }\AttributeTok{{-}s}\NormalTok{ /run/fcgi.sock /usr/bin/fcgiwrap }\KeywordTok{\&\&} \DataTypeTok{\textbackslash{}}
    \ExtensionTok{nginx} \AttributeTok{{-}g} \StringTok{"daemon off;"}
\end{Highlighting}
\end{Shaded}

\textbf{startup.sh}

\begin{Shaded}
\begin{Highlighting}[]
\CommentTok{\#!/bin/sh}

\FunctionTok{chmod} \AttributeTok{{-}R}\NormalTok{ 777 /git}
\BuiltInTok{cd}\NormalTok{ /git}
\FunctionTok{git}\NormalTok{ init }\AttributeTok{{-}{-}bare} \VariableTok{$\{1\}}\NormalTok{.git}
\BuiltInTok{cd} \VariableTok{$\{1\}}\NormalTok{.git}
\FunctionTok{git}\NormalTok{ branch }\AttributeTok{{-}m}\NormalTok{ main}
\FunctionTok{git}\NormalTok{ config http.receivepack true }

\BuiltInTok{exit}\NormalTok{ 0}
\end{Highlighting}
\end{Shaded}

Der Container wirs über folgenden Befehl erstellt und gestartet:

\begin{Shaded}
\begin{Highlighting}[]
\ExtensionTok{docker{-}compose}\NormalTok{ up }\AttributeTok{{-}d}
\end{Highlighting}
\end{Shaded}

Das Skript \emph{startup.sh} erstellt im Container ein neues Repository.
Es kann entweder vom Hostsystem oder aus dem Container gestartet werden:

\textbf{Hostsystem}

\begin{Shaded}
\begin{Highlighting}[]
\ExtensionTok{docker}\NormalTok{ exec }\AttributeTok{{-}it}\NormalTok{ git{-}nginx{-}container /startup.sh projekt  }\CommentTok{\# Reponame OHNE .git angeben}
\end{Highlighting}
\end{Shaded}

\textbf{Im Container}

\begin{Shaded}
\begin{Highlighting}[]
\ExtensionTok{docker}\NormalTok{ exec }\AttributeTok{{-}it}\NormalTok{ git{-}nginx{-}container /bin/sh}
\ExtensionTok{./startup.sh}\NormalTok{ projekt  }\CommentTok{\# OHNE Endung .git}
\end{Highlighting}
\end{Shaded}

\emph{Verlassen des Containers}

\begin{Shaded}
\begin{Highlighting}[]
\ExtensionTok{strg}\NormalTok{ + pq}
\end{Highlighting}
\end{Shaded}

\textbf{Clonen}\\
Zum Clonen wird dann verwendet

\begin{Shaded}
\begin{Highlighting}[]
\FunctionTok{git}\NormalTok{ clone http://}\OperatorTok{\textless{}}\NormalTok{IP}\OperatorTok{\textgreater{}}\NormalTok{:8080/git/projekt.git}
\end{Highlighting}
\end{Shaded}





\end{document}
